% Plantilla de presentación - PFM02 / IIIEPE
% Alumno: Mtro. Martin Jonathan de la Cruz Muñoz
% Matrícula: 00874
%
% ==============================================================================
% PLANTILLA BASE DE ACTIVIDADES - PFM02
% Temas selectos de matemáticas I
% Maestría en Aprendizaje y Enseñanza de las Matemáticas
% Marzo-julio 2026, Grupo A
% ==============================================================================
% Uso recomendado:
% 1. Duplica este archivo por actividad y conserva el original como plantilla.
%    Ejemplo: PFM02_S03_Enteros_SecuenciaDidactica.tex
% 2. Actualiza el título visible de portada, la sección de actividad editable y la
%    fecha de entrega cuando corresponda.
% 3. Conserva las diapositivas de identidad PFM02 cuando la entrega requiera
%    contextualizar la unidad; si no son necesarias, comenta la sección completa.
% 4. No pegues la consigna completa en el producto final. Usa la consigna como
%    lista de cumplimiento y redacta una respuesta académica propia.
% 5. Toda actividad debe cerrar con evidencia matemática y reflexión didáctica.
% 6. Si la actividad pide secuencia, ficha, tabla, mecanizaciones, situación en
%    contexto o instrumento de evaluación, constrúyelo en LaTeX/TikZ cuando sea viable.
%
% Enfoque académico para PFM02:
% - Partir de un contenido matemático escolar y su proceso de desarrollo de aprendizaje.
% - Relacionar concepto, procedimiento, representación, situación en contexto y evaluación.
% - Evitar hojas de ejercicios sin intención didáctica: antecedente -> explicación ->
%   práctica -> situación contextualizada -> evaluación -> metacognición.
% - Integrar el Programa Sintético Fase 6 y el campo Saberes y Pensamiento Científico
%   cuando la secuencia se vincule con secundaria.
% - Priorizar una voz docente: precisa en lo matemático, clara en lo didáctico y
%   cuidadosa en la progresión de dificultad.
%
% Estructura sugerida por actividad:
% [ ] Portada con sesión, tema, alumno, matrícula, grupo y fecha.
% [ ] Contenido, PDA y aprendizaje de la secuencia.
% [ ] Fase 1: conocimientos previos o antecedente del contenido.
% [ ] Fase 2: exploración, introducción y explicación del contenido.
% [ ] Fase 3: práctica del contenido: mecanizaciones y situaciones en contexto.
% [ ] Fase 4: evaluación: examen, laboratorio, exposición, actividad lúdica,
%     autoevaluación o coevaluación.
% [ ] Fase 5: retroalimentación y reflexión del desempeño/contenido.
% [ ] Fuentes consultadas en formato consistente.
%
% Checklist antes de entregar:
% [ ] Título, actividad, bloque, sesión y fecha están actualizados.
% [ ] El objetivo coincide con la consigna del EVA.
% [ ] El contenido, PDA y aprendizaje esperado están completos.
% [ ] La secuencia respeta las cinco fases de trabajo.
% [ ] Hay ejercicios graduados, situaciones en contexto y cierre metacognitivo.
% [ ] La evaluación corresponde al aprendizaje de la secuencia.
% [ ] Las operaciones, fórmulas, unidades y resultados matemáticos fueron revisados.
% [ ] Ortografía, acentos, nombres propios y fechas revisadas.
% [ ] El PDF compila sin errores y las imágenes institucionales cargan.
% ==============================================================================

\pdfminorversion=7
\documentclass[
	spanish,
	aspectratio=1610,
	xcolor={dvipsnames,table}
]{beamer}

% Lienzo 16:10 ampliado: da más espacio por diapositiva sin cambiar el diseño.
\geometry{paperwidth=19.2cm,paperheight=12cm}

% Idioma, tipografía y recursos visuales
\usepackage[utf8]{inputenc}
\usepackage[T1]{fontenc}
\usepackage[spanish,es-nodecimaldot]{babel}
\usepackage{lmodern}
\usepackage{booktabs}
\usepackage{graphicx}
\usepackage{ragged2e}
\usepackage{tikz}
\usepackage{hyperref}

% Datos del alumno y del programa
\newcommand{\studentname}{Mtro. Martin Jonathan de la Cruz Muñoz}
\newcommand{\studentshort}{M. de la Cruz Muñoz}
\newcommand{\studentid}{00874}
\newcommand{\studentemail}{martin.delacruz@campus.iiiepe.edu.mx}
\newcommand{\programname}{Maestría en Aprendizaje y Enseñanza de las Matemáticas}
\newcommand{\programshort}{MAEM}
\newcommand{\modality}{Modalidad presencial}
\newcommand{\periodname}{1er tetramestre marzo--julio 2026}
\newcommand{\coursecode}{PFM02}
\newcommand{\coursename}{Temas selectos de matemáticas I}
\newcommand{\coursegroup}{Grupo A}
\newcommand{\courseperiod}{Marzo--julio 2026}
\newcommand{\learningunit}{Unidad de aprendizaje}
\newcommand{\evaname}{EVA posgrado}
\newcommand{\blockname}{Secuencia didáctica}
\newcommand{\sessionname}{Sesiones 01--03}
\newcommand{\methodology}{Aprendizaje basado en indagación}
\newcommand{\steamenfoque}{STEAM como enfoque}
\newcommand{\curriculumsource}{Programa Sintético Fase 6}
\newcommand{\formativefield}{Saberes y Pensamiento Científico}
\newcommand{\mainproduct}{Secuencia didáctica matemática}

% Datos institucionales
\newcommand{\institutionname}{IIIEPE}
\newcommand{\institutionlongname}{Instituto de Investigación, Innovación y Estudios de Posgrado para la Educación del Estado de Nuevo León}
\newcommand{\institutionurl}{https://iiiepe.edu.mx/}
\newcommand{\institutionplace}{Monterrey, Nuevo León}

% Recursos oficiales descargados desde el sitio del IIIEPE
\newcommand{\iiiepelogo}{img/iiiepe/logo-iiiepe-blanco.png}
\newcommand{\iiiepemark}{img/iiiepe/marca-iiiepe.png}
\newcommand{\iiiepecover}{img/iiiepe/portada-institucional.png}
\newcommand{\iiiepedots}{img/iiiepe/logo-puntos.png}
\newcommand{\maemplanimage}{img/iiiepe/maem-plan-estudios.png}
\newcommand{\maemsource}{https://iiiepe.edu.mx/maestria-en-aprendizaje-y-ensenanza-de-las-matematicas/}

% Paleta inspirada en la identidad visual del sitio institucional
\definecolor{iiiepeNavy}{HTML}{163B5C}
\definecolor{iiiepeBlue}{HTML}{1F6F9F}
\definecolor{iiiepeAqua}{HTML}{20B7A8}
\definecolor{iiiepeOrange}{HTML}{F28C28}
\definecolor{iiiepePaper}{HTML}{F6F8F7}
\definecolor{iiiepeInk}{HTML}{1E2A32}
\definecolor{iiiepeMuted}{HTML}{667085}

\mode<presentation>{
	\usetheme{default}
	\usefonttheme{professionalfonts}
	\setbeamertemplate{navigation symbols}{}
	\setbeamertemplate{blocks}[rounded][shadow=false]
	\setbeamercovered{invisible}
}

\hypersetup{
	colorlinks=true,
	urlcolor=iiiepeBlue,
	linkcolor=iiiepeNavy
}

\setbeamercolor{background canvas}{bg=white}
\setbeamercolor{normal text}{fg=iiiepeInk,bg=white}
\setbeamercolor{structure}{fg=iiiepeBlue}
\setbeamercolor{frametitle}{fg=white,bg=iiiepeNavy}
\setbeamercolor{title}{fg=white,bg=iiiepeNavy}
\setbeamercolor{subtitle}{fg=white}
\setbeamercolor{block title}{fg=white,bg=iiiepeBlue}
\setbeamercolor{block body}{fg=iiiepeInk,bg=iiiepePaper}
\setbeamercolor{alerted text}{fg=iiiepeOrange}
\setbeamerfont{title}{size=\LARGE,series=\bfseries}
\setbeamerfont{frametitle}{size=\large,series=\bfseries}
\setbeamerfont{block title}{series=\bfseries}
\setbeamertemplate{itemize item}{\textcolor{iiiepeAqua}{\large$\blacktriangleright$}}
\setbeamertemplate{itemize subitem}{\textcolor{iiiepeOrange}{\scriptsize$\blacksquare$}}
\setbeamertemplate{enumerate item}{\textcolor{iiiepeBlue}{\insertenumlabel.}}

\setbeamertemplate{frametitle}{
	\nointerlineskip
	\begin{beamercolorbox}[wd=\textwidth,ht=1.05cm,dp=0.22cm,leftskip=0.35cm,rightskip=0.25cm]{frametitle}
		\begin{minipage}[c]{0.72\textwidth}
			\insertframetitle
		\end{minipage}
		\hfill
		\raisebox{-0.1cm}{\includegraphics[height=0.58cm]{\iiiepelogo}}
	\end{beamercolorbox}
	\vspace{-0.04cm}
	{\color{iiiepeOrange}\rule{\textwidth}{1.4pt}}
}

\setbeamertemplate{footline}{
	\leavevmode%
	\hbox{%
		\begin{beamercolorbox}[wd=.34\paperwidth,ht=2.7ex,dp=1ex,leftskip=1em]{author in head/foot}%
			\color{white}\studentshort
		\end{beamercolorbox}%
		\begin{beamercolorbox}[wd=.38\paperwidth,ht=2.7ex,dp=1ex,center]{title in head/foot}%
			\color{white}\programshort\ · \modality
		\end{beamercolorbox}%
		\begin{beamercolorbox}[wd=.20\paperwidth,ht=2.7ex,dp=1ex,rightskip=1em plus 1fil]{date in head/foot}%
			\color{white}Matrícula \studentid\hfill\insertframenumber/\inserttotalframenumber
		\end{beamercolorbox}%
	}%
	\vskip0pt%
}
\setbeamercolor{author in head/foot}{bg=iiiepeNavy}
\setbeamercolor{title in head/foot}{bg=iiiepeBlue}
\setbeamercolor{date in head/foot}{bg=iiiepeNavy}

\newcommand{\accentbar}{%
	\begin{tikzpicture}[remember picture,overlay]
		\path[use as bounding box] (0,0) rectangle (0,0);
		\fill[iiiepeAqua] (current page.south west) rectangle ([xshift=0.42\paperwidth,yshift=0.08cm]current page.south west);
		\fill[iiiepeOrange] ([xshift=0.42\paperwidth]current page.south west) rectangle ([xshift=\paperwidth,yshift=0.08cm]current page.south west);
	\end{tikzpicture}
}

\newcommand{\sectionframe}[1]{%
	\begin{frame}[plain]
		\begin{tikzpicture}[remember picture,overlay]
			\fill[iiiepeNavy] (current page.south west) rectangle (current page.north east);
			\node[opacity=0.13,anchor=east] at ([xshift=0.8cm]current page.east) {\includegraphics[height=9.4cm]{\iiiepedots}};
			\fill[iiiepeAqua] ([xshift=0.85cm,yshift=2.1cm]current page.south west) rectangle ([xshift=1.05cm,yshift=6.3cm]current page.south west);
			\fill[iiiepeOrange] ([xshift=1.15cm,yshift=2.1cm]current page.south west) rectangle ([xshift=1.35cm,yshift=5.35cm]current page.south west);
		\end{tikzpicture}
		\vspace{1.8cm}
		\hspace{1.25cm}
		\begin{minipage}{0.72\paperwidth}
			{\color{white}\Huge\bfseries #1}\\[0.25cm]
			{\color{white!75}\programname}
		\end{minipage}
	\end{frame}
}

\AtBeginSection[]{
	\sectionframe{\insertsectionhead}
}

\title[\coursecode]{\coursename}
\subtitle{\programname}
\author[\studentshort]{\texorpdfstring{\studentname\\\footnotesize Matrícula: \studentid\\\href{mailto:\studentemail}{\studentemail}}{\studentname, Matricula: \studentid}}
\institute[\institutionname]{\institutionlongname\\\coursecode\ · \coursegroup}
\date[\courseperiod]{\institutionplace\\\courseperiod}

\begin{document}

% Portada institucional
\begin{frame}[plain]
	\begin{tikzpicture}[remember picture,overlay]
		\node[anchor=center,opacity=0.18] at (current page.center) {\includegraphics[width=\paperwidth,height=\paperheight]{\iiiepecover}};
		\fill[iiiepeNavy,opacity=0.94] (current page.south west) rectangle ([xshift=0.55\paperwidth]current page.north west);
		\fill[iiiepeBlue,opacity=0.92] ([xshift=0.55\paperwidth]current page.south west) rectangle (current page.north east);
		\node[anchor=north west] at ([xshift=0.55cm,yshift=-0.45cm]current page.north west) {\includegraphics[width=4.15cm]{\iiiepelogo}};
		\node[opacity=0.17,anchor=south east] at ([xshift=0.9cm,yshift=-0.3cm]current page.south east) {\includegraphics[height=8.8cm]{\iiiepedots}};
		\fill[iiiepeOrange] ([xshift=0.56\paperwidth,yshift=1.15cm]current page.south west) rectangle ([xshift=0.92\paperwidth,yshift=1.28cm]current page.south west);
		\fill[iiiepeAqua] ([xshift=0.56\paperwidth,yshift=0.88cm]current page.south west) rectangle ([xshift=0.82\paperwidth,yshift=1.01cm]current page.south west);
	\end{tikzpicture}

	\vspace{1.9cm}
	\begin{minipage}{0.50\paperwidth}
		{\color{white}\Huge\bfseries Temas selectos\\ de matemáticas I}\\[0.25cm]
		{\color{white!86}\large \coursecode\ · \coursegroup}\\[0.35cm]
		{\color{iiiepeOrange}\rule{0.82\linewidth}{1.3pt}}\\[0.45cm]
		{\color{white}\studentname}\\
		{\color{white!80}\footnotesize Matrícula: \studentid\ · \studentemail}
	\end{minipage}
	\hfill
	\begin{minipage}{0.36\paperwidth}
		\raggedleft
		{\color{white}\Large\bfseries \institutionname}\\[0.18cm]
		{\color{white!82}\footnotesize \institutionlongname}\\[0.5cm]
		{\color{white}\small \programname}\\
		{\color{white!78}\small \courseperiod}\\[0.5cm]
		{\color{white!74}\footnotesize \institutionplace\ · \learningunit}\\
		{\color{white!74}\footnotesize \href{\institutionurl}{\institutionurl}}
	\end{minipage}
\end{frame}

\begin{frame}{Contenido}
	\tableofcontents
\end{frame}

\section{Identidad PFM02}

\begin{frame}{Ficha de la unidad}
	\begin{columns}[T]
		\begin{column}{0.42\textwidth}
			\begin{center}
				\includegraphics[width=0.54\linewidth]{\iiiepemark}\\[0.35cm]
				{\Large\bfseries\color{iiiepeNavy}\institutionname}\\
				{\small\color{iiiepeMuted}\evaname}
			\end{center}
		\end{column}
		\begin{column}{0.54\textwidth}
			\begin{block}{Alumno}
				\textbf{\studentname}\\
				Matrícula: \studentid\\
				Correo: \href{mailto:\studentemail}{\studentemail}
			\end{block}
			\begin{block}{Unidad de aprendizaje}
				\textbf{\coursecode\ · \coursename}\\
				\coursegroup\\
				\courseperiod
			\end{block}
		\end{column}
	\end{columns}
	\accentbar
\end{frame}

\begin{frame}{Contexto académico}
	\begin{columns}[T]
		\begin{column}{0.48\textwidth}
			\begin{block}{Datos generales}
				\begin{itemize}
					\item Programa: \programname.
					\item Unidad: \coursename.
					\item Clave: \coursecode.
					\item Periodo: \courseperiod.
					\item Grupo: \coursegroup.
				\end{itemize}
			\end{block}
		\end{column}
		\begin{column}{0.48\textwidth}
			\begin{block}{Propósito de la unidad}
				Diseñar, analizar y presentar secuencias didácticas para contenidos
				matemáticos de secundaria, articulando contenido, PDA, práctica,
				situaciones en contexto, evaluación y retroalimentación.
			\end{block}
		\end{column}
	\end{columns}
	\accentbar
\end{frame}

\begin{frame}{Enfoque metodológico}
	\begin{block}{Metodología de referencia}
		\methodology\ con \steamenfoque: organizar el aprendizaje desde preguntas,
		exploración, uso de saberes previos, práctica matemática, aplicación y
		reflexión sobre el desempeño.
	\end{block}
	\begin{columns}[T]
		\begin{column}{0.32\textwidth}
			\begin{block}{Inicio}
				Antecedente del contenido, conocimientos previos y preguntas de
				indagación.
			\end{block}
		\end{column}
		\begin{column}{0.32\textwidth}
			\begin{block}{Desarrollo}
				Exploración, explicación, mecanización y situaciones en contexto.
			\end{block}
		\end{column}
		\begin{column}{0.32\textwidth}
			\begin{block}{Cierre}
				Evaluación, coevaluación, autoevaluación y metacognición.
			\end{block}
		\end{column}
	\end{columns}
	\accentbar
\end{frame}

\begin{frame}{Base curricular}
	\begin{columns}[T]
		\begin{column}{0.48\textwidth}
			\begin{block}{Documento base}
				\curriculumsource\ para educación secundaria, especialmente el
				campo formativo \formativefield.
			\end{block}
		\end{column}
		\begin{column}{0.48\textwidth}
			\begin{block}{Criterio de uso}
				El Programa Sintético funciona como punto de partida; la secuencia
				debe contextualizar contenido, PDA, evaluación y actividades a las
				necesidades del grupo.
			\end{block}
		\end{column}
	\end{columns}
	\accentbar
\end{frame}

\begin{frame}{Mapa de materiales revisados}
	\begin{enumerate}
		\item Presentación inicial: metodología de indagación y fases de trabajo.
		\item Plantilla DOCX: estructura de secuencia didáctica en cinco fases.
		\item Sesión 01: conversión de fracciones a decimales y viceversa.
		\item Sesión 02: suma y resta con fracciones.
		\item Sesión 03A: multiplicación y división con fracciones.
		\item Sesión 03B: operaciones con enteros.
		\item Sesión 11A: volumen de prismas y pirámides.
	\end{enumerate}
	\accentbar
\end{frame}

\section{Secuencias revisadas}

\begin{frame}{Sesión 01}
	\begin{block}{Conversión de fracciones a decimales y viceversa}
		Contenido: expresión de fracciones como decimales y de decimales como
		fracciones.
	\end{block}
	\begin{columns}[T]
		\begin{column}{0.48\textwidth}
			\begin{block}{PDA}
				Usa diversas estrategias al convertir números fraccionarios a
				decimales y viceversa.
			\end{block}
		\end{column}
		\begin{column}{0.48\textwidth}
			\begin{block}{Aprendizaje}
				Convertir fracciones a decimales y decimales a fracciones.
			\end{block}
		\end{column}
	\end{columns}
	\accentbar
\end{frame}

\begin{frame}{Sesión 02}
	\begin{block}{Suma y resta con fracciones}
		Contenido: extensión del significado de las operaciones y sus relaciones
		inversas.
	\end{block}
	\begin{columns}[T]
		\begin{column}{0.48\textwidth}
			\begin{block}{PDA}
				Reconoce el significado de las cuatro operaciones básicas y sus
				relaciones inversas al resolver problemas.
			\end{block}
		\end{column}
		\begin{column}{0.48\textwidth}
			\begin{block}{Aprendizaje}
				Resolver situaciones problemáticas que involucran sumar o restar
				fracciones.
			\end{block}
		\end{column}
	\end{columns}
	\accentbar
\end{frame}

\begin{frame}{Sesión 03A}
	\begin{block}{Multiplicación y división con fracciones}
		Contenido: extensión del significado de las operaciones y sus relaciones
		inversas.
	\end{block}
	\begin{itemize}
		\item Aprendizaje: resolver situaciones problemáticas que involucran
		multiplicar o dividir fracciones.
		\item Práctica: operaciones directas, interpretación del procedimiento y
		revisión de resultados.
		\item Cuidado didáctico: evitar que ``invertir y multiplicar'' quede como
		regla aislada sin significado.
	\end{itemize}
	\accentbar
\end{frame}

\begin{frame}{Sesión 03B}
	\begin{block}{Operaciones con enteros}
		Contenido: extensión de los números a positivos y negativos y su orden.
	\end{block}
	\begin{itemize}
		\item PDA: reconoce la necesidad de los números negativos usando cantidades
		que tienen al cero como referencia.
		\item Aprendizaje: resolver sumas, restas, multiplicaciones y divisiones de
		números con signo.
		\item Situaciones en contexto: balance financiero, profundidad bajo el nivel
		del mar y variaciones respecto a cero.
	\end{itemize}
	\accentbar
\end{frame}

\begin{frame}{Sesión 11A}
	\begin{block}{Volumen de prismas y pirámides}
		Contenido: medición y cálculo en diferentes contextos.
	\end{block}
	\begin{columns}[T]
		\begin{column}{0.48\textwidth}
			\begin{block}{PDA}
				Usa diferentes estrategias para calcular el volumen de prismas,
				pirámides y cilindros.
			\end{block}
		\end{column}
		\begin{column}{0.48\textwidth}
			\begin{block}{Aprendizaje}
				Calcular el volumen de prismas y pirámides mediante fórmulas,
				unidades y problemas contextualizados.
			\end{block}
		\end{column}
	\end{columns}
	\accentbar
\end{frame}

\section{Modelo de secuencia}

\begin{frame}{Estructura de secuencia didáctica}
	\begin{block}{Datos de entrada}
		Tema, contenido, PDA, aprendizaje de la secuencia, nombre del estudiante,
		grupo y fecha.
	\end{block}
	\scriptsize
	\begin{tabular}{p{0.20\textwidth}p{0.70\textwidth}}
		\toprule
		\textbf{Elemento} & \textbf{Función didáctica} \\
		\midrule
		Contenido & Precisa el objeto matemático escolar que se abordará. \\
		PDA & Describe el proceso esperado de desarrollo de aprendizaje. \\
		Aprendizaje & Traduce el PDA en logro operativo de la secuencia. \\
		Evidencia & Muestra cómo se verificará comprensión, procedimiento y uso. \\
		\bottomrule
	\end{tabular}
	\accentbar
\end{frame}

\begin{frame}{Cinco fases de trabajo}
	\begin{columns}[T]
		\begin{column}{0.48\textwidth}
			\begin{block}{Fase 1}
				Conocimientos previos: recuperar antecedentes, errores frecuentes y
				preguntas iniciales.
			\end{block}
			\begin{block}{Fase 2}
				Explicación del contenido: explorar, representar, formalizar y
				ejemplificar.
			\end{block}
		\end{column}
		\begin{column}{0.48\textwidth}
			\begin{block}{Fase 3}
				Práctica: resolver mecanizaciones y situaciones en contexto.
			\end{block}
			\begin{block}{Fases 4 y 5}
				Evaluar, retroalimentar y promover reflexión del desempeño personal.
			\end{block}
		\end{column}
	\end{columns}
	\accentbar
\end{frame}

\begin{frame}{Matriz editable de secuencia}
	\scriptsize
	\begin{tabular}{p{0.18\textwidth}p{0.30\textwidth}p{0.26\textwidth}p{0.16\textwidth}}
		\toprule
		\textbf{Fase} & \textbf{Actividad docente} & \textbf{Actividad del estudiante} & \textbf{Evidencia} \\
		\midrule
		1. Previos & Plantea problema inicial & Explica estrategias conocidas & Respuestas iniciales \\
		2. Explicación & Modela procedimiento & Registra pasos y ejemplos & Apuntes y ejercicios \\
		3. Práctica & Propone retos graduados & Resuelve y justifica & Procedimientos \\
		4. Evaluación & Aplica instrumento & Demuestra desempeño & Producto/rúbrica \\
		5. Reflexión & Retroalimenta & Reconoce logros y dudas & Metacognición \\
		\bottomrule
	\end{tabular}
	\accentbar
\end{frame}

\begin{frame}{Criterios matemáticos}
	\begin{itemize}
		\item Precisar el significado del concepto antes de mecanizar.
		\item Usar representaciones diversas: simbólica, verbal, gráfica, tabular o geométrica.
		\item Graduar dificultad: ejemplo guiado, práctica directa, variación y problema contextual.
		\item Revisar unidades, signos, equivalencias, simplificación y sentido del resultado.
		\item Anticipar errores frecuentes y convertirlos en retroalimentación.
	\end{itemize}
	\accentbar
\end{frame}

\section{Plantilla editable}

\begin{frame}{Portada de actividad}
	\begin{block}{Sustituir datos}
		Tema: escribir aquí el contenido matemático de la sesión.
	\end{block}
	\begin{columns}[T]
		\begin{column}{0.48\textwidth}
			\begin{block}{Identificación}
				\coursecode\ · \coursename\\
				\coursegroup\ · \courseperiod\\
				Alumno: \studentshort
			\end{block}
		\end{column}
		\begin{column}{0.48\textwidth}
			\begin{block}{Producto}
				\mainproduct\\
				Sesión: escribir número\\
				Fecha: escribir fecha
			\end{block}
		\end{column}
	\end{columns}
	\accentbar
\end{frame}

\begin{frame}{Contenido, PDA y aprendizaje}
	\begin{block}{Contenido}
		Escribir el contenido matemático tal como se trabajará en la secuencia.
	\end{block}
	\begin{block}{Proceso de desarrollo de aprendizaje}
		Escribir el PDA correspondiente, cuidando que sea observable y coherente
		con el Programa Sintético o con la consigna del curso.
	\end{block}
	\begin{block}{Aprendizaje de la secuencia}
		Redactar qué será capaz de resolver, explicar, justificar o aplicar el
		estudiante al terminar.
	\end{block}
	\accentbar
\end{frame}

\begin{frame}{Fase 1: conocimientos previos}
	\begin{block}{Propósito}
		Identificar antecedentes del contenido, estrategias espontáneas y posibles
		errores antes de presentar el procedimiento formal.
	\end{block}
	\begin{itemize}
		\item Pregunta detonadora o situación inicial.
		\item Registro de saberes previos.
		\item Diagnóstico rápido: ejercicio, discusión, ejemplo o predicción.
	\end{itemize}
	\accentbar
\end{frame}

\begin{frame}{Fase 2: explicación del contenido}
	\begin{block}{Propósito}
		Introducir el contenido con representaciones, lenguaje matemático y
		ejemplos graduados.
	\end{block}
	\begin{itemize}
		\item Definición o regla matemática.
		\item Ejemplo guiado paso a paso.
		\item Contraejemplo o error frecuente.
		\item Pregunta de verificación conceptual.
	\end{itemize}
	\accentbar
\end{frame}

\begin{frame}{Fase 3: práctica del contenido}
	\begin{block}{Propósito}
		Resolver mecanizaciones y situaciones en contexto para pasar del
		procedimiento a la interpretación.
	\end{block}
	\begin{itemize}
		\item Ejercicios directos para consolidar procedimiento.
		\item Problemas con contexto escolar, financiero, geométrico o cotidiano.
		\item Justificación del método usado.
		\item Comparación de estrategias cuando sea posible.
	\end{itemize}
	\accentbar
\end{frame}

\begin{frame}{Fase 4: evaluación}
	\begin{columns}[T]
		\begin{column}{0.48\textwidth}
			\begin{block}{Instrumentos posibles}
				Examen breve, laboratorio, exposición, actividad lúdica,
				autoevaluación, coevaluación o rúbrica.
			\end{block}
		\end{column}
		\begin{column}{0.48\textwidth}
			\begin{block}{Evidencia esperada}
				Procedimiento correcto, argumentación, interpretación del resultado
				y transferencia a una situación nueva.
			\end{block}
		\end{column}
	\end{columns}
	\accentbar
\end{frame}

\begin{frame}{Fase 5: retroalimentación}
	\begin{block}{Reflexión del desempeño}
		Usar preguntas de cierre para que el estudiante identifique qué aprendió,
		qué estrategia le funcionó, qué dificultad permanece y cómo resolvería un
		problema similar.
	\end{block}
	\begin{itemize}
		\item Hoy aprendí que...
		\item El paso más difícil fue...
		\item Una estrategia que me ayudó fue...
		\item Si volviera a resolver un problema parecido, primero haría...
	\end{itemize}
	\accentbar
\end{frame}

\section{Pautas de entrega}

\begin{frame}{Ruta sugerida para crear una secuencia}
	\begin{enumerate}
		\item Elegir contenido y PDA.
		\item Definir aprendizaje de la secuencia con verbo observable.
		\item Diseñar fase 1 con saberes previos y pregunta detonadora.
		\item Explicar el contenido con ejemplos y representaciones.
		\item Graduar práctica: mecanización, variación y contexto.
		\item Evaluar y cerrar con metacognición.
	\end{enumerate}
	\accentbar
\end{frame}

\begin{frame}{Checklist PFM02}
	\begin{itemize}
		\item Tema, contenido, PDA y aprendizaje están alineados.
		\item Las cinco fases aparecen completas.
		\item Los ejercicios tienen progresión de dificultad.
		\item Hay al menos una situación en contexto.
		\item La evaluación mide el aprendizaje declarado.
		\item La retroalimentación permite reflexión del contenido y desempeño.
	\end{itemize}
	\accentbar
\end{frame}

\begin{frame}{Producto solicitado}
	\begin{columns}[T]
		\begin{column}{0.58\textwidth}
			\begin{itemize}
				\item Insertar aquí la secuencia didáctica o síntesis visual.
				\item Incluir fórmulas, operaciones, tablas o problemas contextualizados.
				\item Cerrar con evaluación y retroalimentación.
			\end{itemize}
		\end{column}
		\begin{column}{0.36\textwidth}
			\begin{block}{Enfoque}
				Contenido matemático + PDA + cinco fases + evidencia de aprendizaje.
			\end{block}
		\end{column}
	\end{columns}
	\accentbar
\end{frame}

\section{Cierre}

\begin{frame}{Conclusiones}
	\begin{enumerate}
		\item Sintetizar el contenido matemático trabajado.
		\item Explicar cómo la secuencia favorece comprensión, práctica y transferencia.
		\item Plantear una mejora didáctica para una nueva aplicación.
	\end{enumerate}
	\accentbar
\end{frame}

\begin{frame}[plain]
	\begin{tikzpicture}[remember picture,overlay]
		\fill[iiiepeNavy] (current page.south west) rectangle (current page.north east);
		\node[opacity=0.13,anchor=south east] at ([xshift=0.8cm,yshift=-0.2cm]current page.south east) {\includegraphics[height=9cm]{\iiiepedots}};
		\node[anchor=north west] at ([xshift=0.55cm,yshift=-0.45cm]current page.north west) {\includegraphics[width=3.7cm]{\iiiepelogo}};
	\end{tikzpicture}
	\vspace{2.2cm}
	\begin{center}
		{\color{white}\Huge\bfseries Gracias}\\[0.55cm]
		{\color{white}\Large \studentname}\\[0.18cm]
		{\color{white!78}\footnotesize Matrícula \studentid\ · \studentemail}\\[0.45cm]
		{\color{iiiepeOrange}\rule{0.36\paperwidth}{1.4pt}}\\[0.35cm]
		{\color{white!72}\footnotesize \institutionname\ · \coursecode\ · \coursename}
	\end{center}
\end{frame}

\end{document}
